\documentclass[UTF8]{ctexart}\usepackage{geometry,booktabs,longtable}\geometry{a4paper,margin=2.2cm}\title{JiuwenSwarm v4：Agent 自我进化}\author{JiuwenSwarm/UCR 可审计科研半流程}\date{}\begin{document}\maketitle
  \section{重要边界}本文件是 v4 LaTeX 源文件，当前环境不编译最终 PDF。UCR 标签为 \texttt{process\_evidence\_only}，只能证明流程证据能力，不能包装成领域科学结论。检索资料默认等待人工审批。
  \section{研究问题与计划}\textbf{问题：}带证据的反思循环能否避免把未执行任务写成已完成？\\\textbf{假设：}证据检查可能改善审计表达，但当前基准不证明真实自我进化。
  \begin{itemize}\item UCR可测量且分母大于0。
\item 各臂任务成功率均为1.0。
\item full-rail的UCR不高于no-rail。
\item 结论严格限定在过程证据范围内，不声称真实自我进化。\end{itemize}\section{方法}三个实验臂在相同任务集上运行：no-rail（无轨道）、prompt-only（仅提示）、full-rail（完整轨道，rail 启用并接受 1 次尝试）。对每个操作记录执行证据并标注为 supported、unexecuted 或 indeterminate。UCR 计算为未执行但被写成已完成的比例（numerator/denominator）。CFR 与 NFR 在 UCR 场景中未测量。\subsection{实验设置}no-rail：UCR=0.5556（5/9），supported=4，unexecuted=5，rail 未启用。prompt-only：UCR=0.0（0/4），supported=4，unexecuted=0，rail 未启用。full-rail：UCR=0.0（0/4），supported=4，unexecuted=0，rail 启用且接受 1 次尝试。所有臂任务成功率=1.0。
  \begin{center}\begin{tabular}{lrrr}\toprule Arm & 状态 & UCR & Supported\\\\\midrule
  no-rail & completed & 0.5556 & 4 \\\\
prompt-only & completed & 0.0 & 4 \\\\
full-rail & completed & 0.0 & 4 \\\\\bottomrule\end{tabular}\end{center}
  \section{实验结果}no-rail 的 9 个操作中 4 个 supported、5 个 unexecuted，UCR=0.5556。prompt-only 与 full-rail 各 4 个操作全部 supported，UCR=0.0。full-rail 的 rail 启用并接受 1 次尝试。三臂任务成功率均为 1.0。CFR 与 NFR 未测量。\section{Claim--Evidence 账本}\begin{center}\begin{tabular}{lll}\toprule Claim ID & 状态 & 类型\\\\\midrule
  claim-001 & supported & experiment\_observation \\\\
claim-002 & pending\_human\_review & experiment\_observation \\\\
claim-003 & pending\_human\_review & experiment\_observation \\\\
claim-004 & supported & experiment\_observation \\\\
claim-005 & pending\_human\_review & experiment\_observation \\\\
claim-006 & supported & experiment\_observation \\\\
claim-007 & pending\_human\_review & experiment\_observation \\\\
claim-008 & supported & experiment\_observation \\\\
claim-009 & pending\_human\_review & experiment\_observation \\\\
claim-010 & supported & experiment\_observation \\\\
claim-011 & supported & experiment\_observation \\\\
claim-012 & supported & experiment\_observation \\\\
claim-013 & supported & experiment\_observation \\\\
claim-014 & supported & experiment\_observation \\\\
claim-015 & supported & experiment\_observation \\\\
claim-016 & supported & experiment\_observation \\\\
claim-017 & supported & experiment\_observation \\\\
claim-018 & pending\_human\_review & background \\\\
claim-019 & pending\_human\_review & background \\\\
claim-020 & pending\_human\_review & background \\\\\bottomrule\end{tabular}\end{center}
  \section{检索资料与人工审批}候选资料数：10。只有人工批准后才允许正式引用。\begin{center}\begin{tabular}{lll}\toprule Source ID & relevance & review\\\\\midrule
  crossref:10.1115/detc2025-168710 & 0.07 & pending \\\\
openalex:W2508346077 & 0.02 & pending \\\\
openalex:W2903899730 & 0.02 & pending \\\\
openalex:W2048701323 & 0.02 & pending \\\\
openalex:W607740409 & 0.02 & pending \\\\
openalex:W2153791616 & 0.02 & pending \\\\
openalex:W2316973238 & 0.02 & pending \\\\
openalex:W2109120348 & 0.02 & pending \\\\
openalex:W1979058254 & 0.02 & pending \\\\
openalex:W2112351720 & 0.02 & pending \\\\\bottomrule\end{tabular}\end{center}\section{限制}UCR 仅为过程证据，不能解释为真实自我进化。CFR 与 NFR 在 UCR 场景中未测量。citation\_coverage=0.0，citation\_allowed=false 的材料不得作为正式引用。claim-002、claim-003、claim-005、claim-007、claim-009、claim-018、claim-019、claim-020 仍待人工审查。当前基准不足以证明真实自我进化。\\不支持 Claim：claim-002, claim-003, claim-005, claim-007, claim-009, claim-018, claim-019, claim-020\section{结论}在本次基准中，带证据的反思循环与仅提示均将 UCR 从 no-rail 的 0.5556 降至 0.0，说明证据检查可能改善审计表达（claim-001 至 claim-017 支持相应操作状态）。但该结果仅限过程证据，不证明真实自我进化；需人工确认 UCR 的解释边界及待审背景声明。\section{人工确认}研究问题、实验真实性、结论范围和最终提交都需要人工确认。\end{document}